\begin{abstract}
  Because a missed defaulter costs several times a false alarm, a default model
  must be judged at an operating point, not by ranking alone. We compare ten
  classifiers on the UCI credit-card default data (29,965 Taiwanese card holders,
  22.1\% defaulters) against a $K$-nearest-neighbour benchmark and a logistic
  baseline, every choice made on training rows and each configuration scored
  once on a frozen 30\% test set. Tree ensembles rank best (XGBoost: test
  AUC 0.789, average precision 0.578) but gain only 0.025 AUC over the logit, the
  random forest is indistinguishable from XGBoost, and six months of repayment
  status carry nearly all the signal, additively. Class weights, SMOTE and
  under-sampling inside the folds move the intercept, not the ranking. At a 5:1
  cost ratio, moving the logit's cut-off from 0.5 to 0.15 saves 1,937 [1,604,
  2,218] of its 7,063 test cost units; switching to XGBoost saves a further 145
  [$-$18, 299].
\end{abstract}

\paragraph{Contribution statement.}
Group B19. Model families are owned end to end (code, tuning, interpretation,
presentation): RUAN Hangjing the KNN benchmark, the logistic family, nested
cross-validation and PCA/K-means; LU Xinyu the pruned tree, random forest and
bagging, gradient boosting and XGBoost, TreeSHAP, monotone constraints and
partial dependence; LI Jingwen the support-vector machines, the neural network,
LDA/QDA/naive Bayes, the class-imbalance study and leakage demonstration, Platt
recalibration, and the freeze and aggregation code. Shares per section (RUAN/LU/LI, \%):
\S1~\textcolor{red}{30/20/50}; \S2~\textcolor{red}{25/15/60};
\S3~\textcolor{red}{30/25/45}; \S4~\textcolor{red}{30/35/35};
\S5~\textcolor{red}{15/60/25}; \S6~\textcolor{red}{40/15/45};
\S7~\textcolor{red}{34/33/33}; code and reproducibility
\textcolor{red}{20/25/55}; presentation \textcolor{red}{33/33/34}
(Table~\ref{tab:a-workload}; red = draft awaiting the group's confirmation). Not shared with any other module.

\section{Introduction}

Unsecured revolving credit is one of the largest sources of loss for a retail
bank: a defaulter's balance is written off, while the upside on a good customer
is a thin monthly margin. The two errors are not symmetric, so a model that merely ranks clients
well is not yet a decision; the question is which model, at which cut-off.

The study is comparative in the sense of the course [6] rather than a hunt for
a winner. Two models are the reference points --- a $K$-nearest-neighbour
classifier, the non-parametric benchmark, and a logistic regression, the
scorecard a credit department could defend to a regulator --- and every other
family is read against them, each gap explained by a mechanism rather than a
rank. Ranking is judged by ROC-AUC and average precision, calibration by the
Brier score and the decision by expected cost; accuracy is context only. Three
methods beyond the syllabus --- TreeSHAP [7], a monotonically constrained
XGBoost and partial dependence [4] --- carry the inference. Four questions
organise what follows. (1)~\emph{Value of complexity}: how much ranking quality do
non-linear and ensemble methods add over the two reference models, beyond
test-set noise? (2)~\emph{Operating point}: cut at 0.5, at maximum $F_1$, or at
the threshold implied by an asymmetric cost --- and does that choice outweigh the
model's? (3)~\emph{Class imbalance}: do class weights, SMOTE and under-sampling
inside the folds help, and what does resampling before the split fabricate?
(4)~\emph{Inference}: which characteristics drive default, in which direction, and
do the coefficient, impurity, permutation and Shapley accounts agree?
